\subsection*{小结}

\begin{table}[htbp]\centering\small
\begin{tabularx}{\linewidth}{@{}lXX@{}}
\toprule
\textbf{概念} & \textbf{语法} & \textbf{作用} \\
\midrule
软件 DMA & \texttt{dma.copy} / \texttt{dma.copy.swiz\textless{}N\textgreater{}} & 线程协作的分块传输 \\
重排 & 拷贝上的 \texttt{.swiz\textless{}N\textgreater{}} 与 \texttt{mma.load.swiz\textless{}N\textgreater{}} & 无布局冲突的共享内存布局；对不同搬运路径效果相同 \\
任意窗口 & \texttt{view(M,N).from(r,c)} & 不规整或运行时定位的切片 \\
带步长的分块 & \texttt{.subspan().step().at()} & 非紧凑布局、相互重叠的模板 \\
不完整分块 & \texttt{.zfill} & 把越界元素补零 \\
形状重新解释 & \texttt{span\_as([dims])} & 暂存缓冲区的零拷贝变形 \\
\bottomrule
\end{tabularx}
\end{table}
